\paragraph{A necessary rank bound from the quadratic terms.}
For a pure kernel direction \(B\), put
\(W_B=\ker T_B/\langle B\rangle\), of dimension five. The coefficients
\(Q_j(X)=[t^2]q_j(B+tX)\), for \(j=2,3,4\), define quadratic forms
on \(W_B\). Indeed, replacing \(X\) by \(X+cB\) adds
\((j-1)c(dq_j)_BX+\binom j2 c^2q_j(B)\), which vanishes on the
polar kernel. A completing three-plane through \(B\) gives a
two-dimensional subspace of \(W_B\) on which all three forms vanish.

Let \(S=k[\lambda_2,\lambda_3,\lambda_4]\), where
\(k=\overline{\mathbb F}_5\), and let
\(\Phi_B=\det(\lambda_2Q_2+\lambda_3Q_3+\lambda_4Q_4)\) in any
basis of \(W_B\), using the symmetric matrices of the forms. Set
\(J_B=(\Phi_B,\partial_{\lambda_2}\Phi_B,
\partial_{\lambda_3}\Phi_B,\partial_{\lambda_4}\Phi_B)\).
We keep the quintic itself because its degree equals the characteristic.
If a completing plane exists, then
\[
                 \dim_k(S/J_B)_9\ \geq\ 3 .
\]
To see this, begin the basis with the two common isotropic directions.
The matrix of the net then has upper-left block zero, upper-right block
\(A\) of size \(2\times3\), and lower-right block \(T\).
Let \(v\) be the column of signed maximal minors of \(A\).
Determinant expansion gives \(\Phi_B=v^{\mathsf T}Tv\).
Consequently \(J_B\subseteq I=(v_1,v_2,v_3)\). If \(v=0\), the
bound follows immediately. Otherwise the two rows of \(A\) give
linear relations on the minors, independent over the fraction field.
Their degree-six polynomial multiples give \(2\dim S_6\) independent
relations among the degree-seven multiples of the three quadrics.
Hence \(\dim I_9\le 3\dim S_7-2\dim S_6=52\), while
\(\dim S_9=55\). The test is an explicit \(78\times55\) matrix of
degree-nine multiples of the four generators of \(J_B\).
A rank of at least 53 excludes a completion over \(k\).

\begin{proposition}[Kernel lines over the quadratic extension]
\label{prop:quadratic-kernel-exclusion}
No geometric level three-plane in the collision branch at five
has a kernel line defined over \(\mathbb F_{25}\). Over \(\mathbb C\)
such a kernel line cannot be defined over a number field of degree
at most two.
\end{proposition}
\begin{proof}
Every nonzero kernel edge vector can be relabeled and scaled so its
\(01\) entry is one. Testing all \(25^5\) assignments of the other
edges against the three pure equations yields 904 normalized directions.
Restoring the projective relabelings gives 964 lines. Under \(S_4\)
and fifth power they form 31 orbits, of sizes \(8,12,24,48\) with
multiplicities \(2,1,17,11\). These operations preserve geometric
extendibility. The quadratic rank test excludes 27 orbits. Its ranks
on the other four are \(52,50,49,48\). For each of these four, use the
ten echelon cells of \(\operatorname{Gr}(2,5)\) to parametrize a line
and impose vanishing of all three quadratic forms on it. Each quadratic
contributes its two diagonal and one mixed coefficient. Exact
Gr\"obner reduction gives \([1]\) in all forty systems, with
\(\omega^2=2\) representing the coefficient extension. Thus the
quadratic conditions alone exclude a common line in all 31 orbits.

For each of the 31 representatives the checker also solves the three
polar rows for a companion with prescribed star \((0,1,0)\). After
substitution the six remaining two-column equations are polynomials
in three free entries. Their multiples through degree four give 39
rows on 35 monomials. The checker computes coefficients for a linear
combination and verifies the identity \(1=\sum h_i f_i\) coefficient by
coefficient over \(\mathbb F_{25}\). Thirty matrices have rank 35.
The remaining matrix has rank 34 and still contains the constant in
its row span. Each identity persists after extension of scalars.

For the complex assertion, choose a place above five of a field of
definition of degree at most two. Its residue field embeds in
\(\mathbb F_{25}\). Extend the kernel line over the local discrete
valuation ring and apply the closed incidence and properness argument
of Proposition~\ref{prop:rational-kernel-exclusion}.
\end{proof}

\begin{proposition}[Dimension of the quadratic direction locus]
\label{prop:kernel-curve}
In \(\mathbb P^5_{\overline{\mathbb F}_5}\), the locus \(Z\) of pure
kernel directions whose three quotient quadrics contain a common line
has dimension at most one. The complex locus defined by the same
pure, polar and quadratic completion conditions also has dimension
at most one.
\end{proposition}
\begin{proof}
The full polar rank is three on the pure locus, by relabeling a nonzero
edge and using the normalized rank calculation above. Hence the spaces
\(W_B\) form a rank-five bundle there. The common-line incidence is
closed in its relative Grassmannian. Its image \(Z\) is closed by
properness.

Put \(l=e_{02}+e_{12}+e_{23}\) and \(m=e_{03}+e_{13}+e_{23}\). Their
common zero set \(\Pi\) is a projective three-space with edge vectors
\[
                 (a,-b-d,-c-d,b,c,d).
\]
Restrict \(q_2,q_3,q_4\) to this space. In graded reverse lexicographic
order \(a>b>c>d\) their initial ideal is
\[
       (ac^5,\ c^6,\ abc^3,\ a^2c^3,\ a^2bc,\ a^3,\ b^2).
\]
There are 24 standard monomials in \(a,b,c\), with counts
\(1,3,5,6,5,3,1\) in degrees zero through six. Since the initial
generators are independent of \(d\), the Hilbert function is constant
at 24 from degree six on. The geometric projective intersection has
length 24. The quadratic-field census supplies 24 distinct points on
\(\Pi\), including the points with \(e_{01}=0\). They exhaust the
intersection and none belongs to \(Z\), by the common-line checks in
Proposition~\ref{prop:quadratic-kernel-exclusion}. Every closed subset
of dimension at least two in \(\mathbb P^5\) meets \(\Pi\)
\cite{stacksProjectiveIntersection}. Hence \(\dim Z\le1\).

For the complex assertion, use the projective flag incidence over
\(\mathbb Z_{(5)}\) of a kernel line \(\langle B\rangle\) in a
tangent three-plane \(L\). Impose the pure equations on \(B\), the
polar and quadratic equations on every vector in \(L\), and \(l=m=0\)
on \(B\). These conditions are closed and are defined at five.
The incidence has empty geometric special fiber by the calculation
above. Properness rules out a complex point, so the projective
dimension argument applies to the complex direction locus as well.
\end{proof}

\paragraph{A finite fiber to test.}
The two linear forms give a morphism \(Z\to\mathbb P^1\) by
\([B]\mapsto[l(B):m(B)]\). Every fiber is projective and contained
in the complement of \(\Pi\) in a hyperplane of the pencil, an affine
four-space. Thus every fiber is finite \cite{stacksproperaffine}.
The fiber \([0:1]\) has the four-coordinate presentation
\[
                 B=(a,-b-d,1-c-d,b,c,d).
\]
The census finds 40 pure directions over \(\mathbb F_{25}\) with this
normalization, all excluded. Theorem~\ref{thm:fixed-face-quadratic-empty}
excludes the entire geometric fiber. Thus \(Z\) is finite, since every
curve component would map onto \(\mathbb P^1\).

\paragraph{An equation on the remaining kernel directions.}
Give the prescribed companion star a homogeneous scale \(h\), so
its entries are \((0,h,0)\). Together with the six free entries of
the companion \(C\) these are seven homogeneous coordinates. In
\(\operatorname{per}(J_4+sB+tC)\) take the coefficients of
\(s^it^j\) for
\[
 (i,j)=(1,1),(2,1),(3,1),(0,2),(1,2),(2,2),(0,3).
\]
Let \(R(B)\) be their projective resultant. An actual companion with
\(h=1\) forces \(R(B)=0\). The companion degrees are
\(1,1,1,2,2,2,3\), and the kernel degree in each block is \(i\).
The multidegree formula gives total kernel degree
\(24(1+2+3)+12(0+1+2)=180\) when the specialized polynomial is
nonzero \cite[Introduction and Lemma~3.4]{dandreaDickensteinResultant}.
At \(B=(3,0,2,1,1,2)\), the three polar forms eliminate three
homogeneous coordinates. The remaining three quadrics and one cubic
generate all 84 sextic monomials in the four remaining coordinates,
as checked by exact row reduction at five. Thus they have no common
projective zero, proving \(R(B)\ne0\).

Every kernel direction of a full plane must lie in \(Z\) and satisfy
all vertex relabelings of \(R\). These resultant equations give a
second way to cut any curve components that survive the quadratic tests.

\paragraph{The pure curve above \([0:1]\).}
Work over \(k=\mathbb F_5\) and write \(m_r=\sum_{i\ne r}e_{ir}\)
with unordered edge indices. Thus \(l=m_2\) and \(m=m_3\).
Before imposing a common line, the closure of the normalized pure locus is
the curve \(C\subset\mathbb P^4\) with edge vector
\[
             (a,-b-d,m-c-d,b,c,d),\qquad q_2=q_3=q_4=0 .
\]
This curve is geometrically integral of degree 24 and arithmetic genus 49.
Indeed, its section at \(m=0\) consists of the 24 reduced points above,
all with \(d\ne0\). Thus \(q_2,q_3,q_4,m,d\) are homogeneous parameters
in a polynomial ring on five variables and form a regular sequence.
In particular \(C\) is a complete intersection. The 24 standard monomials
of the section at \(m=d=0\) lift to a basis of its homogeneous coordinate
ring over \(k[m,d]\), by graded generation and the complete-intersection
Hilbert series. Hence \([m:d]\) gives a finite flat projection of degree
24. The counts \((1,3,5,6,5,3,1)\) in degrees zero through six give
Hilbert polynomial \(24n-48\).

Two other projective hyperplane sections have the following closed points.
All occur with multiplicity one.
\[
\begin{array}{c|c|c}
 \text{hyperplane} & \text{residue degrees over } k
                    & \text{degree-nine Jacobian ranks}\\ \hline
 d=4m & 1,1,2,20 & 54,54,54,54\\
 d=a  & 6,18     & 54,54
\end{array}
\]
In each case the lexicographic basis on \(m=1\) expresses \(a,b\) as
polynomials in \(c\) and has a squarefree degree-24 eliminant in \(c\).
The second hyperplane has no point at infinity, as substitution of
\(a=d=1,m=0\) gives the unit ideal. For the first this follows from
\(d\ne0\) at infinity. A reduced finite hyperplane section meets every
curve component in smooth points of the curve. Thus \(C\) is generically
reduced and is reduced by the complete-intersection property. The
degree-20 point lies on a unique component over \(k\), of degree at
least 20. Any remaining components would have total degree at most four,
whereas their section at \(d=a\) would contain a point of degree at
least six. Hence \(C\) is irreducible over \(k\). Its smooth rational
point lies on a unique geometric component, fixed by Frobenius.
Transitivity on geometric components and perfectness of \(k\) give
geometric integrality. The two disjoint sections exclude 48 geometric
directions from \(Z\). In the chart \(m=1\), the Jacobian ideal of \(C\)
has length ten and support \(a=c=0,\ d=1,\ b^2+b+1=0\).
The reduced section at infinity gives smoothness there.

\paragraph{The smooth model and its Jacobian.}
The normalization \(\nu:\widetilde C\to C\) has genus 43. Each of the
two geometric singularities has two smooth branches with intersection
multiplicity three. The branches at either point are already defined
over its residue field \(K=\mathbb F_{25}\).

To see this, write \(s^2+s+1=0\) and work at
\((a,b,c,d)=(0,s,0,1)\) in the chart \(m=1\). Set \(t=d-1\), and let
\[
 M=\det\frac{\partial(q_2,q_4)}{\partial(a,b)},\qquad
 G=\det\frac{\partial(q_2,q_4,q_3)}{\partial(a,b,c)}.
\]
At this point \(M=-2s-1\ne0\). The formal implicit function theorem
solves \(q_2=q_4=0\) for \(a,b-s\) in \(K[[c,t]]\). Write \(F(c,t)\)
for \(q_3\) on this smooth surface. The chain rule gives \(G=M F_c\)
on the surface. The determinant
\(\det\bigl(\partial(q_2,q_4,G)/\partial(a,b,c)\bigr)\) at the point
is \(-1\). In particular \(F_{cc}(0,0)=-1/M^2=2\), and there is a
unique formal critical arc \(c=c_0(t)\) with \(F_c=0\), \(c_0(0)=0\).

The following rows give the coefficients of \(t,\ldots,t^6\) on this
arc. Direct substitution in the original permanent equations gives
\(q_2=q_4=G=0\bmod t^7\) and \(q_3=-2t^6\bmod t^7\).
\[
\begin{array}{c|rrrrrr}
 & t&t^2&t^3&t^4&t^5&t^6\\ \hline
 a   &1&1&0&2&0&-2\\
 b-s &2&s-2&-2s-1&2-s&2-s&2-s\\
 c   &s&2-s&2-s&2s+1&-2s-1&s-2
\end{array}
\]
Thus after translation by \(c_0(t)\) the equation is
\(X^2U(X,t)-t^6 V(t)=0\) with \(U(0,0)=1\) and \(V(0)=2\).
The unit \(2U/V\) has a square root with constant term one, since
2 is invertible. The new coordinate \(Y=X\sqrt{2U/V}\) gives
\[
        \widehat{\mathcal O}_{C_K,(0,s,0,1)}
                   \simeq K[[Y,t]]/(Y^2-2t^6).
\]
Indeed \((2s+1)^2=2\). Evaluation along the branches
\(Y=\pm(2s+1)t^3\) identifies this ring with the pairs of series in
\(K[[t]]\times K[[t]]\) congruent modulo \(t^3\). Its normalization
is the whole product, and the quotient has length three. Completion
preserves this delta invariant \cite{stacksDeltaCompletion}. Frobenius
gives the same calculation at the conjugate singularity. By the
normalization exact sequence and the previously computed arithmetic
genus, \(g(\widetilde C)=49-3-3=43\). One can construct this model over \(\mathbb F_5\) by blowing up the
reduced singular locus three times, taking the new reduced singular
locus as center at each step. In the completed charts put
\(Y=tY_1,\ Y_1=tY_2,\ Y_2=tY_3\) successively. The strict transforms
satisfy \(Y_1^2=2t^4,\ Y_2^2=2t^2,\ Y_3^2=2\). The first two transforms
each have a unique singular point above the center. At the third step
the two tangent directions are finite and distinct. The centers are
Frobenius-stable and the morphisms are isomorphisms off those centers.
The resulting smooth projective geometrically connected curve is the
normalization.

The Jacobian has dimension 43. For each prime \(\ell\ne5\) its Tate
module has rank 86 \cite{stacksPicardCurves}. The computed point counts
on \(C\) are 18 over \(\mathbb F_5\) and 64 over \(\mathbb F_{25}\).
The two singular points and their four branch points are all defined
over \(\mathbb F_{25}\), with none rational over \(\mathbb F_5\).
The counts on \(\widetilde C\) are consequently 18 and 66. If \(\pi\)
is the fifth-power Frobenius on its Jacobian, the curve trace formula
gives \(\operatorname{Tr}(\pi)=-12\) and
\(\operatorname{Tr}(\pi^2)=-40\). Its zeta numerator
\(\det(1-T\pi\mid V_\ell\operatorname{Jac}(\widetilde C))\) has degree
86 and begins \(1+12T+92T^2\) \cite{longCurveFrobenius}.
For an ample line bundle \(L\) on the projective Jacobian over
\(\mathbb F_5\), let \(H\) be the numerical class of the symmetric
ample bundle \(L\otimes[-1]^*L\). Frobenius pulls back \(H\) to \(5H\),
so \(\pi^\dagger\pi=[5]\) for its Rosati adjoint. For rational \(a,b\)
the trace identity \(2gQ_H(u)=\operatorname{Tr}(u^\dagger u)\) and
\(\operatorname{Tr}(\pi^\dagger)=\operatorname{Tr}(\pi)\) give
\[
 Q_H(a\,1+b\pi)=a^2-\frac{12}{43}ab+5b^2,\qquad
 Q_H(\pi-1)=\frac{270}{43}.
\]
This computes the positive degree form on the span of the identity
and Frobenius \cite{longAbelianComparison}. Its discriminant is
\((12/43)^2-20<0\), so two rational linear combinations of these
endomorphisms agree exactly when their two coefficients agree. Applying its equality
test to the common-line condition still requires constructing two
endomorphisms and proving that their equality detects that incidence.


\paragraph{What the two sections decompose.}
Let \(T\) be either of these finite sections of pure directions and write
\(A=\mathcal O(T)=k[c]/(h)\) for its squarefree shape eliminant.
The residue degrees give isomorphisms of \(k\)-algebras
\[
 \mathcal O(T_{d=4m})\simeq
 k\times k\times\mathbb F_{5^2}\times\mathbb F_{5^{20}},
 \qquad
 \mathcal O(T_{d=a})\simeq\mathbb F_{5^6}\times\mathbb F_{5^{18}}.
\]
For each irreducible factor \(f\mid h\), the class of
\((h/f)((h/f)^{-1}\bmod f)\) projects onto that field factor. These
idempotents are orthogonal and sum to one. Each factor corresponds to
one Frobenius orbit of geometric points. In particular the second
section has 24 geometric points and no \(k\)-rational point.

There is also an intrinsic diagonal decomposition for the map
\(T\to\operatorname{Spec} k\). Its self-product has coordinate ring
\(D=k[x,y]/(h(x),h(y))\), with diagonal ideal \(J=(x-y)\).
The divided difference
\[
    e_\Delta=\frac{h(x)-h(y)}{x-y}\,h'(x)^{-1}\quad\text{in }D
\]
is defined since \(h\) is squarefree. It restricts to one on the
diagonal and vanishes on all distinct pairs of roots. After splitting
\(h\), these evaluations account for the entire ring. Descent gives
\(J=(1-e_\Delta)\) and \(\operatorname{Ann}(J)=(e_\Delta)\). Hence
\(J\cap\operatorname{Ann}(J)=0\) and \(J+\operatorname{Ann}(J)=D\),
as in the Galois obstruction criterion \cite{gotCollisionDefect}.
The residual algebra \(D/(e_\Delta)\) represents the 552 ordered
distinct pairs over \(\bar k\), alongside 24 diagonal pairs.

These products describe the two finite sections. The ambient pure
curve is integral. The geometric common-line incidence introduces the
further requirement that a direction support a two-space for the three
quadrics. Its equivalence with the affine chart below concerns existence
on this fixed face. Turning a decomposition of a witness family into
a polynomial-time selector would additionally require a uniform algorithm
that chooses and recovers a witness from each encoded yes-instance.

\paragraph{Testing lines over the residue fields.}
Consider the pencil \(d=a+rm\). Its finite sections have no point at
infinity, since \(m=0,\ d=a\) misses \(C\). At \(r=0\) the two residue
fields above account for all 24 geometric directions. In particular,
over \(\mathbb F_{5^e}\) the number of pure directions in this section is
\[
   6\,\mathbf{1}_{6\mid e}+18\,\mathbf{1}_{18\mid e}.
\]
To test for a common line, take one representative of each closed point
and form its three quotient quadrics over the residue field. On both
factors the polar sextic described below has an explicit inverse
modulo the three hyperplane quadrics. The checker multiplies the inverse
by the sextic and reduces the product to one. The same identity holds
over every field extension. Thus even over \(\mathbb F_{5^{18}}\),
where all 24 directions are rational, none admits a geometric common
line. This gives a second certificate for the rank exclusions at \(r=0\).

The same test excludes the four translated sections. The table lists
the degree of each closed point over \(\mathbb F_5\). All have geometric
multiplicity one except for the two marked quadratic points.
\[
\begin{array}{c|l|r}
 r & \text{residue degrees} & \text{geometric points} \\ \hline
 0 & 6,18 &24\\
 1 & 2\ (\mathrm{mult.}\ 2),\ 2\ (\mathrm{mult.}\ 4),\ 6,6 &16\\
 2 & 3,3,4,4,5,5 &24\\
 3 & 1,1,1,1,1,1,2,2,7,7 &24\\
 4 & 2,22 &24
\end{array}
\]
For \(r=1,2,3\), eliminate first to \(b+c\), then separate the points
by \(b\) when needed. For \(r=4\), \(c\) separates the points. Factoring
the eliminants gives the residue fields, and standard monomials give
the lengths of the corresponding factors. The line test uses the
residue field at each point, including those with multiplicity.
The five sections are disjoint, so their total length 120 accounts for
112 distinct geometric directions. The fiber at \(m=0\) was excluded
above. Consequently every common-line direction on this face satisfies
\[
           m\ne0,\qquad ((d-a)/m)^5-(d-a)/m\ne0.
\]

To test recovery when a line exists, consider the three quadrics
\(x_1^2-2x_0^2,\ x_3^2-2x_2^2,\ x_4^2\) on \(\mathbb P^4\) over
\(\mathbb F_5\). For \(\alpha^2=2\), their geometric support consists of
four lines obtained by choosing the signs in \(x_1=\pm\alpha x_0\)
and \(x_3=\pm\alpha x_2\), with \(x_4=0\). Frobenius exchanges both
signs, so there is no line defined over \(\mathbb F_5\). Over
\(\mathbb F_{25}\) the search returns a frame for one line and checks
its rank and all nine Gram pairings. The search uses the ten row-echelon
charts for two-planes. Over a specified finite field it can exhaust
the coordinates of each chart. A maximal ideal containing a proper
chart ideal has residue field finite over \(\mathbb F_5\), so searching
successive extension degrees divisible by the coefficient-field degree
eventually recovers a frame. The uniform selector target still requires
a polynomial bound in the encoded input length.

\paragraph{Edge conditions on a surviving direction.}
Every \(B\in Z\) with \(l(B)=0\) has all six edges nonzero. The
permutations fixing vertex 2 have two edge orbits, represented by
\(e_{12}\) and \(e_{13}\). Relabeling preserves the projective face,
and the excluded pencil base section permits normalization at \(m=1\).
When \(e_{12}=0\), the pure equations give closed points of degrees
8 and 12. When \(e_{13}=0\), the (degree, multiplicity) pairs are
\((1,2),(1,3),(2,2),(2,3),(4,1)\). The separating coordinates for
these lexicographic calculations are \(d\) and \(d+2b\), respectively.

Here is the common-line test at each of these seven closed points. In
the quotient space of dimension five, a common line meets the hyperplane
\(x_4=0\). For the chosen frames, on \(x_3=1\) the three restricted
quadrics reduce all monomials of degree four to terms of degree at most
three. Homogenization also excludes any point at \(x_3=0\) in that
hyperplane. At a common zero \(x\), form the signed cofactor vector
\(y\) of the four rows \(Q_2x,Q_3x,Q_4x,e_3^*\).
If the first three rows have rank three, their kernel must equal
any common two-space through \(x\), so \(Q_2(y)=0\). If their rank is
smaller, \(y=0\). In each residue field, multiples of the reduced sextic
\(Q_2(y)\) span all eight remaining monomials modulo the three quadrics.
This gives a unit and excludes the line over the algebraic closure.

There is a second reduction under the same vertex permutations. Let
\(x_i=e_{2i}\), for \(i=0,1,3\), and let \(y_i\) be the edge opposite
\(x_i\). The \(x_i\) sum to zero. If two coincide, then either all
vanish (the pure equations give a unit), or the equal pair can be
scaled to \(1,1\). The latter section has four rational points and two
closed points of degree nine, all of Jacobian rank 54 in degree nine.
For distinct \(x_i\), write \(y_i=u+v x_i+w x_i^2\) by interpolation.
For \(w=0,u=0\), the pure equations force the \(x_i\) to be scaled
cube roots of unity with \(v=0\) or \(3\). For \(w=0,u\ne0\),
normalize \(u=1\). The eliminant for \(v\) is irreducible of degree
seven, and the cubic of incident edges splits over the degree-14
extension. Each case gives rank 54, allowing vertex permutations and
Frobenius. Thus a surviving direction has distinct \(x_i\) and \(w\ne0\).
Scaling all edges by \(w\) sets the quadratic coefficient to one.
With \(x_i\) the roots of \(T^3+pT-q\), the three pure equations
now use just \(p,q,u,v\), on \(-4p^3-27q^2\ne0\). For instance the
quadratic equation becomes
\[
 (4v+3)q+2p^2+2pu+2pv^2+3pv+2p+u^2=0 .
\]
Interpolating the incident, opposite and outer diagonal entries of a
tangent vector in the same cubic and setting the incident linear
coefficient to zero gives eight coordinates modulo \(B\). Reduction in
the ordered splitting algebra gives three linear polars and three
quadratic forms in these eight coordinates, with coefficients in
\(k[p,q,u,v]\).

On the surviving locus \(q\ne0\), since it is the product of the incident
edges. With \(u=0\) the pure quotient gives two closed points of degrees
5 and 9, both excluded by rank 54. With \(v=3\) the degrees are 2, 2
and 7. One quadratic point has a zero opposite edge, and the other
points again have rank 54. We may therefore solve the displayed
equation for \(q\). Substitute in the two remaining pure equations and
clear powers of \(4v+3\). Their resultant in \(p\), after division by
\(u^4(v-3)^8\), has bidegree \((12,24)\) in \((u,v)\). Any surviving
direction projects to this plane curve.

The resulting polynomial is absolutely irreducible. Its leading
coefficient in \(u\) is a nonzero constant. At \(v=2\) its factors
have degrees 2 and 10, so a proper factorization over \(k[v]\)
would have a degree-two factor in \(u\). At a root of \(v^2+v+1\),
the norm of the specialization to \(k[u]\) factors into irreducibles
of degrees 10 and 14. A degree-two factor would give a degree-four
factor of this norm. Finally the rational point \((u,v)=(4,4)\)
has gradient \((4,2)\), so it is smooth and singles out a unique
geometric component. The first subresultant in \(p\) is
\((v-3)^5(A(u,v)p+D(u,v))\), with bidegrees \((9,19)\) and \((10,17)\)
for \(A,D\). Since \(\deg_u A<12\), it is nonzero in the function
field of the plane curve. Generically one therefore recovers
\(p=-D/A\), then \(q\) from the quadratic equation. On the open \(u\ne0,v\ne3\), a lift with \(A=0\) must also
satisfy \(D=0\). The resultant of \(A\) and the plane equation in
\(u\) has \(v\)-factors of degrees \(1,1,5,9,44\).
At \(v=0\) and at the factors of degrees 5 and 9, the common zeros
of \(A,D\) and the plane equation have \(u=0\). At the degree-44
factor their monic gcd in \(u\) is linear with nonzero constant.
The remaining two equations in \(p\) have an irreducible quadratic
gcd there, verified by its nonsquare discriminant. Thus the exception
is a single closed invariant point of degree 88. The powers
\(1,p,\ldots,p^{87}\) give a basis of its residue field. In that basis
the degree-nine Jacobian rank of the descended net is 54, excluding
every conjugate. Hence \(A\ne0\) and the two rational recovery formulas
apply at every surviving common-line direction on this face.

In particular one can test every prime-field slope \(v\) while allowing
\(u\) to range over the algebraic closure. For \(v=0,1,2,4\), respectively,
the irreducible factors left after the edge and discriminant checks have
degrees \((8)\), \((2,4,4)\), \((2,10)\), and \((2,2,3,4)\).
All ten closed points have degree-nine rank 54. Together with the
\(v=3\) calculation, this excludes every \(v\in\mathbb F_5\).

\paragraph{The quotient by vertex permutations.}
Let \(S_3\) permute the vertices \(0,1,3\), fixing vertex 2. This action
preserves the pure curve \(C\), since it permutes rows and columns of
the permanent simultaneously and preserves the vertex-2 sum.
It lifts to the normalization \(\widetilde C\).
The smooth projective model of the invariant plane curve in \((u,v)\)
is \(\widetilde C/S_3\) and has genus seven.

Indeed, work on the dense open with distinct incident edges, normalized
\(w=1\), and \(qu(v-3)\ne0\). The excluded edge and interpolation
sections above are finite. The projective direction there is recovered
from \(p,q,u,v\) by ordering the three roots of \(X^3+pX-q\) and
setting \(y_i=x_i^2+vx_i+u\). After eliminating \(q\), the first
remaining equation is quartic in \(p\) with leading coefficient 4.
Thus the map to \((u,v)\) has finite fibers and dense image on the
irreducible plane curve. Its first subresultant recovers \(p\)
generically, then the quadratic equation recovers \(q\). The splitting
field of the incident cubic therefore gives \(k(C)\) over \(k(u,v)\),
of degree at most six. All six permutations act faithfully on the
ordered roots and fix \(u,v\), proving \(k(C)^{S_3}=k(u,v)\).

For the ramification, work over \(\bar k\). Use incident and opposite
triples \(x=(x_0,x_1,x_3)\) and \(y=(y_0,y_1,y_3)\). A projective
fixed point of a transposition must be in one of its two eigenspaces.
For \((01)\), the incident triples scale respectively to \((1,1,-2)\)
and \((1,-1,0)\), since the all-zero incident triple is excluded.
The corresponding opposite triples are \((h,h,j)\) and \((h,-h,0)\).
In both cases the three pure equations generate the unit ideal.
For a three-cycle, the invariant eigenvalue would force \(x=0\).
The other two eigenspaces have \(x=(1,\omega,\omega^2)\),
\(y=vx\), for the two primitive cube roots \(\omega\).
Substitution gives \(\omega^2+\omega+1=0\) and \(v^2(v+2)=0\).
The two points with \(v=0\) are the singularities computed above.
The three-cycle fixes each of their two branches, since its order is
three. The two points with \(v=3\) are smooth. This gives six fixed
points on \(\widetilde C\), each with stabilizer of order three, and
three geometric branch values on the quotient. All ramification is
tame in characteristic five. Riemann--Hurwitz \cite{stacksTameHurwitz} gives
\[
       2\cdot43-2=6(2g-2)+6(3-1),\qquad g=7.
\]
All inertia lies in \(A_3\), so the intermediate double cover
\(\widetilde C/A_3\to\widetilde C/S_3\) is everywhere unramified
and has source genus 13. Its function field adjoins the square root of
\(-4p^3-27q^2\), since the Vandermonde product in the ordered roots
changes sign under each odd permutation and squares to the discriminant.

\paragraph{All three relabelings of the pencil.}
Write \(Y=\sum_i y_i\), \(z_i=x_i-y_i\). The three distinct pencils
are \([z_i:m_i]\) with \(m_i=Y+z_i\). The original pencil has
\(i=3\). Every pencil excludes 136 geometric directions on \(C\)
when its section at infinity is included. At \(Y=0\) all pencil values
are 1, giving the 16 geometric points of \(d=a+m\).
Elsewhere normalize \(Y=1\). Then \(\sum_i z_i=-1\), and \([z_i:1+z_i]\)
is rational over \(\mathbb F_5\) exactly when \(z_i\) is rational.
Two rational entries force all three to be rational. For the 25
ordered triples in \(\mathbb F_5^3\) of sum \(-1\), substitution of
\(y_i=x_i-z_i\) in the pure equations gives geometric points only
for permutations of \((1,1,2)\) and \((1,4,4)\), respectively two and
six per ordering. The triple intersection thus has \(16+6+18=40\)
distinct points, and every pairwise intersection equals it. The three
pencils together exclude \(3\cdot136-2\cdot40=328\) geometric
directions. These counts concern points on the original pure curve.
For the finite sections alone, the pattern \((1,4,4)\) contributes
two values at infinity. The finite triple overlap has 22 points and
the finite union has \(3\cdot112-2\cdot22=292\).

This exclusion can be evaluated on the quotient. In the \(w=1\) chart
put \(Y=3u-2p\) and
\[
 H(T)=\operatorname{Res}_X\!\left(X^3+pX-q,\,
             T+X^2-(1-v)X+u\right)=\prod_i(T-z_i).
\]
Every surviving direction satisfies \(Y\ne0\) and
\(H(\lambda Y)\ne0\) for all \(\lambda\in\mathbb F_5\).
This condition tests the three pencils using the invariant coordinates.

Vertex relabeling and Frobenius also act on pairs \((B,L)\) of a
direction and a common line. Projection to \(B\) commutes with both
actions, giving bijections of the line fibers over conjugate directions.
On the free labeling locus, if \(F^e B=gB\) with \(g\in S_3\),
transport around that quotient orbit sends a line over \(B\) to
\(g^{-1}F^eL\). A choice returns to the same line exactly when this
map fixes its two-dimensional subspace. The line equations yield the
global exclusion in Theorem~\ref{thm:fixed-face-quadratic-empty}.


\paragraph{One affine chart for the common-line question.}
Let \(L\) be the three-space above a common line in \(W_B\).
For \(r=0,1,3\), the excluded base section and vertex symmetry imply
\(m_r(B)\ne0\). Thus \(U_r=L\cap\ker m_r\) represents the common
two-space modulo \(B\). For \(r=3\) use the eight coordinates
\[
       X_{01},X_{02},X_{12},X_{13},X_{23},X_{11},X_{22},X_{33}.
\]
Here \(X_{03}=-X_{13}-X_{23}\) and the entry sum determines \(X_{00}\).
Projection to the first three coordinates is injective on \(U_3\).
Indeed, for a nonzero vector with these three entries zero, normalize
its first nonzero entry to one. All five resulting systems, imposing
the pure equations and the three polar and three quadratic equations
on that vector, have reduced Gröbner basis \([1]\).

Suppose no relabeling fixing vertex 2 makes \((X_{01},X_{02})\)
independent on \(U_3\). Swapping 0 and 1 shows that both pairs involving
\(X_{01}\) in the first three coordinates are dependent. The projection
has rank two, so \(X_{01}=0\) on \(U_3\). Applying the same argument with
\(r=0,1\) gives proportionality of the opposite edge and \(m_r\) on \(L\).
Take the unique \(X\in U_3\) with \(X_{02}=0,X_{12}=1\). Then
\[
    m_0(B)X_{13}-B_{13}m_0(X)=
    m_1(B)X_{03}-B_{03}m_1(X)=0 .
\]
Adding these two relations to its pure, polar and quadratic system gives
reduced basis \([1]\) as well. Consequently every common line admits,
after relabeling, a frame \(X,Y\in U_3\) whose \((01,02)\) entries are
\((1,0),(0,1)\). The converse holds by coefficient extraction. The one
remaining system has four kernel coordinates, twelve free frame entries
and eighteen equations obtained from three pure coefficients, six
linear polars and nine quadratic coefficients of \(\operatorname{per}(J_4+sB+tX+zY)\).
Thus emptiness of the zero set of this single ideal over the algebraic closure decides
the quadratic common-line fiber at \([0:1]\). Its unit status remains open.

\paragraph{Finite divisors containing the residual directions.}
There is a degree bound even for the quadratic conditions alone. On
\(C\) the three polar rows have full rank and define an exact sequence
\[
 0\longrightarrow E\longrightarrow\mathcal O_C^9\longrightarrow
 \mathcal O_C(1)\oplus\mathcal O_C(2)\oplus\mathcal O_C(3)
 \longrightarrow 0 .
\]
The kernel direction is the tautological line \(\mathcal O_C(-1)\)
inside \(E\). Its quotient \(W\) therefore has determinant
\(\mathcal O_C(-5)\). The quadratic form \(Q_j\) on \(W\) takes
values in \(\mathcal O_C(j-2)\). Give \(\lambda_2,\lambda_3,\lambda_4\)
weights \(0,1,2\). In the determinant quintic the coefficient of
\(\lambda^\alpha\) is a section of
\((\det W^*)^2\otimes\mathcal O_C(\mathrm{wt}(\alpha))
\simeq\mathcal O_C(10+\mathrm{wt}(\alpha))\).

At \(B=(1,2,2,4,0,4)\) on \(C\), the degree-nine Jacobian matrix has
a nonzero \(53\)-minor with twist 755. To reproduce the choice,
order the columns by increasing weight. Order the rows by increasing
\(10+\mathrm{wt}(\lambda_j)-\mathrm{wt}(\lambda^s)\) for a monomial
multiple \(\lambda^s\partial_j\Phi_B\), omitting the derivative weight
for a multiple of \(\Phi_B\). Keep the first 53 independent rows, then
the first 53 independent columns on those rows, using the enumerations
in the checker for ties. Each determinant term has the same twist,
the sum of the selected row and column weights. This minor is a nonzero
section of \(\mathcal O_C(755)\). Its zero divisor has length
\(24\cdot755=18{,}120\) and contains every direction of the quadratic
collision fiber by its rank ceiling. At the relabeling \((1,0,4,2,2,4)\) of the same point the companion
sextic matrix has rank 84. The corresponding degree-180 resultant thus
cuts out a divisor of length \(24\cdot180=4{,}320\) containing every
direction admitting a full completion \cite{stacksCurveDegrees}.
These lengths count zeros of individual equations with multiplicity.
The surviving directions have distinct incident edges after normalizing
the quadratic interpolation coefficient to one. Their vertex-permutation
orbits on this face therefore have size six, giving at most 3,020 orbits
of quadratic collision directions.
Since \(\mathcal O_C(1)\) is the pullback of \(\mathcal O(1)\) under
\([m:d]\), the norm of the minor is a nonzero binary form of degree
18,120 \cite{stacksFiniteNorm}. Computing it and testing its factors against the remaining
chart would turn the bound into a list of surviving directions.

\paragraph{A finite divisor on the vertex quotient.}
The following equation on the genus-seven quotient gives a finite
index set for testing the remaining common-line directions. On the
rank-five bundle \(W\) above, take the two vertex-invariant linear forms
\[
 H(X)=X_{02}+X_{12}+X_{23},\qquad D(X)=X_{00}+X_{11}+X_{33}.
\]
Both kill \(B\). Indeed \(H(B)=m_2(B)=0\), and if \(T\) is the sum of its
six edges, the corrected diagonal entries are \(T/2-2m_i(B)\).
The sum of the three outer diagonal entries is \(3T/2-4T=0\) at five.
Thus the two forms are defined on \(W_B\).

Write \(b_j\) for the symmetric bilinear form with
\(b_j(x,x)=Q_j(x)\). Wedge the four covectors
\(b_2(x,-), b_3(x,-), b_4(x,-), D\) to get a cofactor vector
\(c(x)\in W\otimes\det W^*\otimes\mathcal O_C(3)\), cubic in \(x\).
Set \(S(x)=Q_2(c(x))\) and form the homogeneous resultant
\[
         R=\operatorname{Res}_x(H,Q_2,Q_3,Q_4,S).
\]
If a two-space \(P\subset W_B\) is common isotropic, take nonzero
\(x\in P\cap\ker H\) over the algebraic closure. The three polar rows
annihilate \(P\). At rank three their common kernel is \(P\), hence
contains the cofactor. At smaller rank the cofactor is zero. In either
case \(S(x)=0\), so \(R(B)=0\). At \(B=(1,1,3,0,3,4)\) the checker uses
\(D,H\) as the final two coordinates of \(W_B\). The quartic reductions
exclude points at infinity in \(H=0\), and the reduced sextic has a
verified inverse on its affine chart. Hence \(R(B)\ne0\).

The five degrees are \((1,2,2,2,6)\), with product 48. The resultant
has degree 48 divided by the corresponding degree in each coefficient
block \cite{dandreaDickensteinResultant}. An invertible change of
variables preserves the universal resultant hypersurface, so it acts on
the resultant by a determinant character. Substitution of a scalar
matrix in the five variables gives exponent 48 on its determinant.
Writing \(M=\det W^*\), the sextic
has coefficient twist \(M^2\otimes\mathcal O_C(6)\). Altogether the
resultant takes values in \(M^{64}\otimes\mathcal O_C(120)\).
Equivariantly \(M=\det(V^*)\otimes\mathcal O_C(5)\), where \(V\) is the
nine-dimensional tangent space before imposing the three polars and
quotienting by \(B\). Its determinant character has square one.
Thus \(R\) is an invariant section of \(\mathcal O_C(440)\).
As \(Y=\sum_i y_i\) is also invariant and nonzero at this test point,
\(h=R/Y^{440}\) is a nonzero rational function on the smooth quotient.
Its pullback has at most \(24\cdot440\) zeros counted with multiplicity.
The degree-six quotient map gives \(\deg (h)_0\le 1760\).
Every common-line direction on the face has \(Y\ne0\).

\begin{theorem}[Decomposition of the common-line incidence on the fixed face]
\label{thm:common-line-decomposition}
Work in characteristic five on the face \([l:m]=[0:1]\).
Let \(U\) be the invariant pure chart with distinct incident edges,
\(w=1\), all six edges nonzero, and \(u(v-3)AY\ne0\), where \(A\)
is the leading coefficient in the first subresultant. It is an open
of the genus-seven quotient with coordinates \(p,q,u,v\). Every
geometric common-line direction on the face lies above \(U\).
Let \(T\to U\) order the three roots of \(X^3+pX-q\).
There is an explicit scheme \(\mathcal F\to U\) of lines in the
descended net. The incidence scheme \(\mathcal I\) of labeled
directions with common lines is the pullback
\[
                \mathcal I\simeq T\times_U\mathcal F .
\]
The projection of \(\mathcal F\) has geometric support in
\(U\cap(h)_0\). At each closed point \(s\) of this candidate divisor
let \(K=\kappa(s)\), let \(e\) be the splitting degree of the incident
cubic over \(K\), and let \(K_e/K\) be the degree-\(e\) extension.
There is an isomorphism of \(K\)-schemes
\[
     \mathcal I_s\simeq
        \coprod_{j=1}^{6/e} (\mathcal F_s\times_K K_e).
\]
In particular \(\mathcal I_s\) is geometrically empty exactly when
\(\mathcal F_s\) is. The fiber \(\mathcal F_s\) retains its scheme structure.
\end{theorem}

\begin{proof}
The exclusions above remove the complement of \(U\) from the possible
incidence. Over its coordinate algebra the ordered-root algebra has
presentation
\[
  x^3+px-q=0,\qquad y^2+xy+x^2+p=0,\qquad z=-x-y .
\]
It is free of rank six by the two monic equations. Inverting the
discriminant makes it an étale \(S_3\)-torsor.
For a tangent vector, evaluate three polynomials of degree at most two
at the ordered roots, one polynomial each for the incident edges, the
opposite edges, and the three outer diagonal entries. The Vandermonde
matrix is invertible. The tangent sum determines the central diagonal.
Adding a multiple of \(B\) shifts the linear coefficient of the
incident polynomial by that multiple. Setting the coefficient to zero
gives eight invariant coordinates. The three polar equations in these
coordinates have constant rank three, since they do after the faithfully
flat extension ordering the roots. Denote their rank-five kernel by
\(E\). In \(\operatorname{Gr}_U(2,E)\) let \(\mathcal F\) be the zero scheme
of the restrictions of the three descended quadrics to the universal
two-space. In a frame these are nine Gram pairings. After pullback to
\(T\) the interpolation is an invertible change of coordinates and
these polynomials are precisely the three original restricted quadrics.
The construction commutes with base change, giving the claimed scheme
isomorphism. The resultant gives the support bound.

Over the residue field, Frobenius permutes the three distinct roots.
On the set of their six orderings it acts freely through the cyclic
group it generates in \(S_3\). All its orbits have size \(e\).
The equivalence between finite étale algebras and finite Galois sets
\cite{stacksEtalePoint} now gives
\[
\begin{array}{c|c|c}
 \text{factor degrees of the cubic over }K & e & \mathcal O(T_s)\\ \hline
 1,1,1 &1&K^6\\
 1,2   &2&K_2^3\\
 3     &3&K_3^2
\end{array}
\]
Pullback gives the formula even for nonreduced or positive-dimensional
line fibers. To retain infinitesimal variation of the direction itself,
let \(Z\subset U\) be the scheme-theoretic image of \(\mathcal F\to U\).
Its support is finite, hence it is a disjoint union of local Artin
schemes. Over each such local ring the same \(6/e\) finite étale
factors lift uniquely up to isomorphism from the residue field
\cite{stacksHenselianEtale}. Pulling back \(\mathcal F\) along those
factors gives the decomposition over the thickened image as well.
\end{proof}

For the degree-\(n\) extension \(K_n/K\), the same formula gives
\(\#\mathcal I_s(K_n)=6\#\mathcal F_s(K_n)\) if \(e\mid n\), and
zero otherwise. Over the algebraic closure every line has six labelings.
The degree bound also gives a finite search independent of an expansion
of \(h\): enumerate the closed points of \(U\) of degree at most 1,760
over \(\mathbb F_5\), recovering \(p,q\) from \(u,v\). Larger-degree points
cannot belong to the support. At each candidate \(s\), the
ten row-echelon charts decide geometric existence in \(\mathcal F_s\).
On a proper chart the field search recovers a line. Over a common
field of definition, evaluating its two rows at the ordered roots
recovers a frame \((B,X_1,X_2)\) and permits direct verification in
the original permanent. If a full constant-permanent plane is required,
impose the remaining cubic and quartic tangent coefficients on the
whole line family. Vanishing of all nonconstant coefficients is
preserved under a change of frame, adding multiples of \(B\) to its
representatives, and permuting the vertices. It therefore defines a
closed subfamily with the same ordering decomposition. For an individual
recovered line these extra coefficients test that particular lift.
The factorization of \((h)_0\) and the list of its nonempty line
fibers remain undetermined.

\paragraph{A complete test over quadratic slopes.}
For each of the ten monic irreducible quadratics \(v^2+sv+t\) over
\(\mathbb F_5\), reducing the invariant plane equation gives
\(a(u)+b(u)v\). In every case \(a,b\) are coprime and
\(a^2-sab+tb^2\) is squarefree of degree 24. Thus \(u\) separates
the geometric points and \(v=-a/b\) recovers the slope. The lifts
satisfy the original pure equations with separable incident cubic.
Across the ten quadratics there are 30 closed quotient points.
Three degree-two points lie on the excluded pencils. At each of the
other 27 the degree-nine Jacobian rank is 54, above the common-line
ceiling of 52. Thus every line fiber over these slopes is geometrically
empty. The ordering algebra has splitting degree one, two and three
at seven, sixteen and seven closed points, respectively. Those give
104 closed labeled directions, comprising 1,440 geometric directions
before testing for a line. With the prime-field slope calculation,
this excludes every \(v\in\mathbb F_{25}\) throughout the eligible chart.


